\section{The Quantum Sector (P7, P17, P31)}
\label{sec:quantum}

\subsection{Quantum statistics from a classical base (P7)}

Bell's theorem says no local model of independent coins matches the quantum
correlations unless the measurement settings are not statistically independent of the
hidden state. A monist mesh with no hidden state denies exactly that independence, since
the settings are made of the same substrate as the system. \textbf{Where we landed,
quantified:} as the setting-state correlation rises, the CHSH value climbs $1.0, 1.6,
2.3, 3.1, 4.0$, crossing the classical bound and passing the quantum value. The currency
is aligned bits, not bits: one bit of aligned dependence reaches $S=4$, one bit
misaligned gives $S=1$. The honest residual: in a natural mesh the shared past of two
spacelike measurements shrinks with separation, so the violation decays with distance
($4.0 \to 2.58 \to 1.04$), while quantum violation does not. The framework makes Bell
violation possible from a classical base and pins the exact remaining puzzle (why the
correlation is aligned and separation-independent), rather than hiding it.

\subsection{Quantum coherence (P17)}

\textbf{Where we landed:} a continuous-time quantum walk on the mesh spreads
ballistically (width $\sim t^{1.00}$) while a classical walk diffuses (width $\sim
t^{0.50}$). Genuine interference, the heart of quantum behaviour, emerges from the
unitary evolution $e^{-iHt}$ on the mesh.

\subsection{The pillars of the formalism (P31)}

\textbf{Where we landed:} the three pillars are present on the mesh. Unitarity, the
total probability $\sum|\psi|^2$ is conserved at $1.000000$. Interference, the
probability of a sum of amplitudes differs from the sum of probabilities by a nonzero
cross term, so amplitudes add, not probabilities. And a conserved Born probability,
$|\psi|^2$ non-negative and summing to one. This is a down-payment, not the full theory.
The open problem is to derive why the probability is $|\psi|^2$ and the measurement rule
from the mesh itself, building on the superdeterministic hinge of P7.
